\documentclass[12pt]{article}
\usepackage[top=1in, bottom=1in, left=1.25in, right=1in]{geometry}
\usepackage{iftex}
\usepackage{graphicx}
\usepackage{enumitem}
\usepackage{xr-hyper}
\usepackage{hyperref}
\externaldocument{appendix4}
\renewcommand{\figureautorefname}{Fig.}
\usepackage[utf8]{inputenc}
\usepackage{newunicodechar}
\newunicodechar{α}{$\alpha$}
\newunicodechar{β}{$\beta$}
\newunicodechar{γ}{$\gamma$}
\usepackage[style=apa, backend=biber]{biblatex}
\addbibresource{fvpollens.bib}
\setcounter{maxnames}{20}
\setcounter{minnames}{1}
\usepackage{color}
\usepackage{amsmath}
\usepackage{amssymb}
\usepackage[export]{adjustbox}
\usepackage{verbatim}
\usepackage{mathpazo}
\usepackage{setspace}
\usepackage{multirow}
\usepackage{lscape}
\usepackage{fancyhdr}
\usepackage[normalem]{ulem}
\usepackage{rotating}
\usepackage{chngcntr}
\usepackage[parfill]{parskip}
\usepackage[tiny,compact]{titlesec}
\usepackage{longtable}
\usepackage{textcomp}
\usepackage{rotating}
\usepackage{xr}
\usepackage{caption}
\usepackage{booktabs}
\usepackage{siunitx}
\usepackage[T1]{fontenc}
\usepackage{gensymb}
\usepackage{float}
\captionsetup{font=footnotesize}
\sisetup{round-mode=places, round-precision=2, detect-all}
\hypersetup{colorlinks=true, linkcolor=black, citecolor=black}
\RequirePackage{lineno}

\def\title{Bumblebees like flowers...presumably?}

\def\author{Jenna B. Melanson$^{1*}$, Natalia Clermont$^{2}$, Tyler T. Kelly$^{2}$, Claire Kremen$^{1,2}$}

\def\keywords{keyword1, keyword2, keyword3}

\def\extras{
  \textbf{Type of article}: Research Article \\
  \textbf{Author contributions}: fill these in \\
  
  \textbf{Acknowledgements:} ...acknowledgements... This research would not have been possible without the support of land owners/managers in the Lower Fraser Valley, along with grant funding from the Natural Sciences and Engineering Research Council of Canada (NSERC) [ALLRP 570736-2021 and RGPIN-2020-05759], the Canada Foundation for Innovation / John R. Evans Leaders Fund (40541) in collaboration with British Columbia Knowledge Development Fund (BCKDF), and the AgriScience Program under Agriculture and Agri-Food Canada's Sustainable Canadian Agricultural Partnership as part of Organic Science Cluster 4 (to CK). JBM was supported by NSERC Vanier CGS 01353 - 000.\\
\textbf{Conflict of Interest:} The authors declare no conflict of interest. \\
\textbf{Data Availability:} Version-controlled code necessary to reproduce the analyses can be found at \url{https://github.com/jmelanson98/fvpollens}. Full datasets will be shared to the same repository and archived on Dryad following acceptance.

}
\clearpage
\def\affiliation{
  \begin{enumerate}
  \item Biodiversity Research Center, University of British Columbia, 2212 Main Mall, Vancouver, BC, Canada, V6T 1Z4
  \item Institute for Resources, Environment, and Sustainability, University of British Columbia, 2202 Main Mall, Vancouver, BC, Canada, V6T 1Z4
  \end{enumerate}
}

\def\runninghead{Bumblebee pollens}

\def\corr{\noindent
  $\star$ Corresponding author: Jenna B. Melanson (jenna.melanson@ubc.ca) \\
  \\
  }

\newcommand{\mstitlepage}{
  \paragraph{Running head:} \textsc{\runninghead}
  \parindent=0pt
  \begin{center}
    {\LARGE \title \par}
    \vskip 3em
    {\large
      \vspace{-20pt}
      \begin{tabular}[t]{c}
        \author
      \end{tabular}\par}
    \vskip 1.5em
  \end{center}\par
  \affiliation
  %\currentaddress
  \corr
  \extras
}

\begin{document}
\mstitlepage
\doublespacing
\linenumbers
\clearpage
\begin{abstract} 
abstract text
\end{abstract} 

\textbf{Keywords:} \keywords
\clearpage
\section{Introduction}
For highly mobile pollinators like bumblebees, colonies acquire food resources (nectars and pollens) within a foraging kernel that can span multiple square kilometres \parencite{melansonAgriculturalIntensificationFavours2026,osborneBumblebeeFlightDistances2008,molaBumbleBeeMovement2025}. In agroecosystems, these foraging kernels encompass a variety of landcover types (e.g., managed crop fields, field or road margins, seminatural habitat fragments) which are likely to provide different floral resources at different times throughout the colony cycle. The importance of landscape-scale floral resource availability for colony growth and reproduction has been demonstrated previously \parencite{williamsLandscapeScaleResourcesPromote2012c,rundlofFlowerPlantingsSupport2022a}, as have the impacts of landscape composition and connectivity on colony-level pollen collection (CITE).

At the level of the individual forager, pollen collection is likely to arise from both local characteristics of the floral assemblage where the forager is captured, as well as landscape-scale habitat availability; importantly, the relative contributions of these two processes may inform our understanding of pollinator behaviour in agricultural landscape mosaics. For example, if foragers show high fidelity to single resource patches (as demonstrated in \textcite{woodgateContinuousRadarTracking2017}), local-scale patterns of floral richness and abundance will best predict the floral richness and turnover of bee-collected pollen loads. In contrast, foragers which link together multiple resource patches in the course of a single foraging trip \parencite{lihoreauRadarTrackingMotionSensitive2012} may collect pollen loads that reflect landscape-scale patterns of floral resource or habitat availability. 

In both cases, floral fidelity is an important secondary consideration; bumble bees have long been thought to ``major and minor'' on floral species, primarily visiting a highly-rewarding focal (``major'') plant species while monitoring the rewards offered by additional non-focal (``minor'') plants \parencite{heinrichBumblebeeEconomics2004}. This behaviour is intuitive from a cognitive standpoint, with multiple lines of evidence suggesting that floral constancy should be the norm (e.g., familiar flowers may provide a stronger search image and allow for more time-efficient extraction of nectar and pollen \parencite{chittkaMindBee2023}). This expectation is supported by empirical evidence highlighting the consistency of mean pollen load species richness across landscapes, timepoints, and bumble bee species in UK farmlands (CITE TIMBERLAKE TENADAY). However, substantial variation in pollen load richness occured around the population mean, suggesting that additional processes---such as the local or landscape-scale availability of the ``major'' plant---may govern floral constancy.

Finally, bumblebees are eusocial species, which collect pollen primarily to provision their young \parencite{goulsonBumblebeesBehaviourEcology2010}. If pollen load species richness is behaviourally constraind at the individual level, as previous evidence suggests, individual pollen load richness (``alpha diversity'') may be an uninformative measure of colony- or population-level intake (``gamma diversity''), which is a proximate determinant of colony growth and reproduction. To better understand the impacts of local and landscape-scale resource availability on pollinator diet breadth, we must also consider the compositional turnover between foragers (``beta diversity'') which may be higher at florally rich sites or in landscapes offering high-quality or diverse habitat types. Not only will this approach provide a clearer understanding of how multiscale resource availability impacts pollinator diet breadth, but will also outline the behavioural processes underpinning this variation, which may tradeoff across space or time.

\section{Methods}
\subsection{Study System}
\subsection{\emph{Bombus} and Floral Surveys}
\subsection{Pollen Metabarcoding}
\subsubsection{Pollen Collection and Quantification}
For all \emph{B. mixtus} and \emph{B. impatiens} workers, we recorded the presence or absence of corbicular pollen as a binary variable. We then removed pollen from the corbiculae of all \emph{B. mixtus} workers and a subset of \emph{B. impatiens} workers; \emph{B. impatiens} workers were selected from the time period which maximized overlap with \emph{B. mixtus} worker activity (18 July-16 August) (\autoref{fig:speciesphenology}). This selection of samples allowed us to draw conclusions about season-long and interannual variation in \emph{B. mixtus} pollen collection, while also allowing us to make comparisons between the two species---which have differing phenologies---during a period of temporal overlap when resource availability was the same for both species. 

We measured the mass of collected pollen loads using a microbalance (GET BRAND) and stored the pollen at -80\degree C until DNA extraction. In cases where only a single pollen load was present (e.g., because the second pollen load was dislodge during bee collection) we multiplied the mass of the single pollen load by two to estimate the total mass of pollen collected by the bee.

\subsubsection{DNA Extraction, Amplification, and Sequencing}



\subsubsection{Bioinformatics}
We processed ITS2 amplicon sequencing data using the DADA2 pipeline \parencite{callahanDADA2HighresolutionSample2016,callahanDADA2ITSPipeline} implemented in R version 4.3.3 \parencite{rcoreteamLanguageEnvironmentStatistical2021a} with libraries \emph{dada2} \parencite{callahanDADA2HighresolutionSample2016}, \emph{Biostrings} \parencite{pagesBiostringsEfficientString2026}, and \emph{ShortRead} \parencite{morganShortReadBioconductorPackage2009}. Given the variable length of ITS sequences across species, we filtered the demultiplexed reads using maxN = 0 but no sequence truncation. We used cutadapt \parencite{martinCutadaptRemovesAdapter2011} to trim primer sequences, then performed a second round of quality filtering using the parameters minLen = 50, maxN= 0, maxEE=c(2,2), truncQ= 2 (DADA2 default parameters \parencite{callahanDADA2ITSPipeline}). The DADA2 pipeline estimates the sequence error rates to probablistically discriminate seqences real sequences from those arising through sequencing errors; we used the learnErrors function to generate a parametric model of sequence errors, which was used to perform inference of the assigned sequence variants (ASVs) using the dada function \parencite{callahanDADA2HighresolutionSample2016}. Finally, we removed chimeric sequences using the removeBimeraDenovo function \parencite{callahanDADA2HighresolutionSample2016}; for all three sequencing runs, the percentage of reads identified as chimeric was between 9-11\%

The output of the DADA2 pipeline is a table of ASVs (i.e., unique sequences) and their abundances across samples. We assigned taxonomy to ASVs in QIIME2 \parencite{bolyenReproducibleInteractiveScalable2019}


\subsection{Statistical Analyses}
Some ideas for analysis:
\begin{itemize}
  \item colony-level bumblebee-floral networks (see \parencite{yeFloralResourcePartitioning2024a}) -- but find some way to account for phenology
\end{itemize}

\section{Results}
\section{Discussion}

\printbibliography
\clearpage


\section*{Acknowledgements}
The authors thank Hazel Barthel and Elinor Sisk? natalia?

%\section*{Conflict of Interest}
%You may be asked to provide a conflict of interest statement during the submission process. Please check the journal's author guidelines for details on what to include in this section. Please ensure you liaise with all co-authors to confirm agreement with the final statement.

\section*{Supporting Information}

\setcounter{figure}{0}
\setcounter{table}{0}


%\input{supplement}

\end{document}
